\section{Small colour classes}\label{sec:recursion}
We first prove Theorem~\ref{thm:gap} when every colour class has at most a small linear number of edges.
This is the case in which the recursion sketched in Section~\ref{sec:sketch-linear} can be carried out on the whole of $K_n$, and in Section~\ref{sec:gap} we will apply it inside small random blocks.

\begin{theorem}\label{thm:linear}
Let $n\ge 2^{128}$. Let $\phi$ be a proper edge-colouring of $K_n$ in which every colour class has at most $2^{-49}n$ edges.
Then $K_n$ contains $\lfloor n/2\rfloor$ pairwise edge-disjoint rainbow spanning trees.
If $n$ is even, they decompose $E(K_n)$.
If $n$ is odd, then, for any prescribed vertex $y$, the trees may be chosen so that the $(n-1)/2$ unused edges are all incident with $y$.
\end{theorem}

\subsection{Reduction to a stronger statement}\label{sec:reduction}
The argument sketched in Section~\ref{sec:sketch-linear} proves the following theorem, which we prove in the rest of this section.
Condition~\ref{Z4} is what allows us to control the quantity $r$ in~\eqref{eq:sketch-r}, while conditions~\ref{Z1}--\ref{Z3} ensure that the number of sets in $\cZ$, and the number of edges with colour in each of them, are small enough for the local lemma.

\begin{theorem}\label{thm:tests}
Let $n\ge2^{128}$ be even. Let $\phi$ be a proper colouring of $K_n$ with every colour class of size at most $M$, where $1\le M\le 2^{-48}n$. Let $\cZ$ be a family of colour sets such that the following hold.
\begin{props}
\item $|\cZ|\le2n$.\label{Z1}
\item $|E_Z(K_n)|\le Mn$ for every $Z\in\cZ$.\label{Z2}
\item Every colour lies in at most $2M$ members of $\cZ$.\label{Z3}
\item For every vertex $v$, there is some $Z_v\in\cZ$ with $C(v)\subseteq Z_v$.\label{Z4}
\end{props}
Let $D=\lceil\sqrt{2Mn}\,\rceil$. Then, $E(K_n)$ decomposes into $n/2$ rainbow spanning trees $T$, each satisfying $\ell_Z(T)\le 2560D+4097(\log n)^2$ for every $Z\in\cZ$.
\end{theorem}

For even $n$, Theorem~\ref{thm:linear} follows by taking $\cZ$ to be the family of sets $C(v)$.
For odd $n$, we remove the prescribed vertex $y$, decompose $E(K_n-y)$ using Theorem~\ref{thm:tests}, and then give each tree one edge at $y$, using the bound on $\ell_{C(y)}(T)$ to find an edge whose colour is new to the tree.

\begin{proof}[Proof of Theorem~\ref{thm:linear} subject to Theorem~\ref{thm:tests}]
For even $n$, apply Theorem~\ref{thm:tests} with $\cZ=\{C(v):v\in V(K_n)\}$. Condition~\ref{Z2} holds because $|E_{C(v)}(K_n)|\le(n-1)M$, condition~\ref{Z3} holds because a colour class covers at most $2M$ vertices, and~\ref{Z1} and~\ref{Z4} are immediate.

For odd $n$, let $y$ be the prescribed vertex and let $m=n-1$, so that $M\le2^{-49}n\le2^{-48}m$.
Apply Theorem~\ref{thm:tests} to $K_n-y$ with $\cZ=\{C(v):v\in V(K_n)\}$, where the sets $C(v)$ are the sets of colours at $v$ in $K_n$.
Then~\ref{Z1}--\ref{Z4} hold with $m$ in place of $n$. Indeed, $n\le 2m$, $|E_{C(v)}(K_n-y)|\le mM$, each colour lies in at most $2M$ of the sets $C(v)$, and the set of colours at $v$ in $K_n-y$ is contained in $C(v)$.
By~\eqref{eq:consts} in Section~\ref{sec:recursion}, every resulting tree $T$ satisfies $\ell_{C(y)}(T)\le L(m)\le m/1024$.
The $m$ edges at $y$ have distinct colours, and, for each tree $T$, at most $m/1024$ of them have colours already used by $T$.
We now take the $m/2$ trees one at a time, and give each tree an edge at $y$ which has not been given to an earlier tree and whose colour does not appear in the tree. Such an edge always exists, since $m-m/1024-m/2>0$.
The resulting $m/2=\lfloor n/2\rfloor$ trees are edge-disjoint rainbow spanning trees of $K_n$, and the unused edges are the remaining $m/2=(n-1)/2$ edges at $y$, as required.
\end{proof}

We now prove Theorem~\ref{thm:tests}, following the sketch in Section~\ref{sec:sketch-linear}.
In Section~\ref{sec:hierarchy}, we fix the shape of the recursion and the parameters which control it.
In Section~\ref{sec:placing}, we place the vertices into the recursion at random, so that every set at the bottom of the recursion spans a rainbow complete graph. We also ensure that no set in the recursion contains many edges of one colour or many edges with colour in some $Z\in\cZ$.
In Section~\ref{sec:leaves}, we decompose the sets at the bottom of the recursion into Hamilton paths, and, in Section~\ref{sec:combining}, we combine these decompositions up the recursion using Lemma~\ref{lem:attach}.

Throughout the rest of this section, let $n$, $\phi$, $M$, $\cZ$ and $D$ be as in Theorem~\ref{thm:tests}, and let $V=V(K_n)$.

\subsection{The shape of the recursion and its parameters}\label{sec:hierarchy}
Let $B$ be the smallest even integer with $B\ge2^{24}(\log n)^2$.
We organise the set $[n]$ into a rooted binary tree $\cT$, whose nodes are intervals in $[n]$, as follows. The root of $\cT$ is $[n]$. Each node of size $N>2B$ has two children, of sizes $2\lfloor N/4\rfloor$ and $N-2\lfloor N/4\rfloor$, which partition it. Nodes of size at most $2B$ have no children, and are the \emph{leaves} of $\cT$.
We record the following simple properties of $\cT$.
All node sizes are even.
A child of a node of size $N>2B$ has size in $[B,\,N/2+1]\subseteq[B,0.6N]$, so that the leaves of $\cT$ have sizes in $[B,2B]$.
Along any path from the root to a leaf of $\cT$, the sizes of the nodes decrease by a factor of at least $5/3$ at each step. Hence, such a path has at most $1+\log_{5/3}n\le2\log n$ nodes, and, on such a path, the sizes of the nodes below a node of size $N$ sum to at most $2.5N$. Finally, $\cT$ has fewer than $n$ nodes.

We now define the four functions which control the recursion. For a node of size $N$, $\Lambda(N)$ will bound the number of edges with colour in any $Z\in \cZ$ inside the corresponding vertex set (see~\eqref{eq:D2}), and $m(N)$ will bound the number of edges of any one colour inside it (see~\eqref{eq:D3}). Further, $t(N)$ will be the value of $t$ when we apply Lemma~\ref{lem:attach} at this node, and $L(N)$ will bound $\ell_Z(T)$ for every tree $T$ we construct on this vertex set and every $Z\in\cZ$ (see~\eqref{eq:load}). Let
\begin{gather*}
\Lambda(N)=16DN^2/n+32N\log n,\qquad m(N)=\lceil32MN^2/n^2+32\log n\rceil,\\
t(N)=\lceil1024DN/n+2048\log n\rceil+1,\qquad L(N)=2560DN/n+4097(\log n)^2.
\end{gather*}
As $n\ge2^{128}$, we have $\log n\ge 88$, and we will use the following bounds:
\begin{equation}\label{eq:consts}
D\le2^{-23}n,\qquad S:=\frac{64MB^3}{n^2}\le\frac{2^{33}(\log n)^6}{n}\le\frac1{16},\qquad
\frac{L(N)}{N}\le2^{-10}\quad\text{and}\quad \frac{t(N)}{N}\le\frac12\quad\text{for all }N\ge B.
\end{equation}
Indeed, as $M\le 2^{-48}n$, we have $\sqrt{2M n}\le \sqrt{2\cdot2^{-48}}\,n=2^{-23.5}n$, and $B\le2^{25}(\log n)^2$. The function $(\log n)^6/n$ is decreasing for $\log n>6$ and is very small when $n= 2^{128}$. Next, $L(N)/N\le2560\cdot2^{-23}+4097/2^{24}<2^{-10}$ for $N\geq B$. Finally, $t(N)/N\le2^{-13}+2048\log n/B+2/B<1/2$ for $N\geq B$.

\subsection{Placing the vertices at random}\label{sec:placing}
Let $\sigma:V\to[n]$ be a uniformly random bijection. For each node $W$ of $\cT$, we use the vertex set $\sigma^{-1}(W)$, which, for simplicity, we also denote by $W$.
We will show that, with positive probability, $K[W]$ is rainbow for every leaf $W$ of $\cT$, and, moreover, that, conditioned on this, every node of $\cT$ also satisfies~\eqref{eq:D2} and~\eqref{eq:D3} below with positive probability.
All the events we consider are canonical events, in the sense of Section~\ref{sec:bijections}, for the single random bijection $\sigma$.

\smallskip

\noindent\textbf{The events to avoid.}
Let $e$ and $e'$ be distinct edges of the same colour, noting that their four vertices are distinct as $\phi$ is proper. Let $f$ be an injection from these four vertices into a single leaf of $\cT$. Let $A_{e,e',f}$ be the event that $\sigma|_{V(e)\cup V(e')}=f$, which is a canonical event with probability $1/\ff n4\le2/n^4$.
Avoiding all of these events means precisely that $K[W]$ is rainbow for every leaf $W$ of $\cT$.

To apply the local lemma, we count the events adjacent to a fixed one. Fix a term $x\mapsto y$. There are at most $nM$ pairs of edges with the same colour such that one of the edges contains $x$, and, letting $b_j$, $j\in [k]$, be the sizes of the leaves of $\cT$, there are at most $\sum_j\ff{b_j}4\le8nB^3$ injections into leaves. So, at most $8Mn^2B^3$ of the events have $x$ in their domain.
There are at most $\sum_\gamma\binom{|E_\gamma|}2\le Mn^2/4$ pairs of edges with the same colour, and at most $4(2B)^3$ injections of four labelled vertices into the leaf containing $y$ that use $y$. So, at most $8Mn^2B^3$ of the events have $y$ in their range.
Let $Q=16Mn^2B^3$, and give every event $A_{e,e',f}$ the weight $x_A=4/n^4$. Each event has at most $4Q$ neighbours, so that $\prod(1-x)\ge1-4Qx_A=1-4S\ge3/4$ by~\eqref{eq:consts}. Since $x_A\cdot\frac34\ge2/n^4$, Lemmas~\ref{lem:LLL} and~\ref{lem:negdep} show that the event $\Omega_0$ that none of the events $A_{e,e',f}$ occurs has positive probability.

\smallskip

\noindent\textbf{Probabilities conditioned on $\Omega_0$.}
Let $E$ be a canonical event with $s\le n/4$ terms. Then $E$ is adjacent to at most $sQ$ of the events $A_{e,e',f}$. By Lemmas~\ref{lem:cond} and~\ref{lem:negdep}, and using that $(1-x)^{-1}\le e^{2x}$ for $x\in [0,1/2]$, we have
\begin{equation}\label{eq:condprob}
\Prb(E\mid\Omega_0)\le\frac{(1-4/n^4)^{-sQ}}{\ff ns}\le\frac{e^{2sS}}{(3n/4)^s}\le\Bigl(\frac2n\Bigr)^s.
\end{equation}
Summing over the at most $\ff Ns$ injections into a node of size $N$, we see that, conditioned on $\Omega_0$, the probability that $s$ given vertices all lie in a given node of size $N$ is at most $(2N/n)^s$.

\smallskip

\noindent\textbf{Edges with colour in a set $Z\in\cZ$.}
For each $Z\in\cZ$, the graph with edge set $E_Z(K_n)$ has at most $Mn$ edges by~\ref{Z2}, and hence has degeneracy at most $\sqrt{2Mn}\leq D$. So, before choosing $\sigma$, we can fix an orientation of each such graph with maximum outdegree at most $D$.
For a vertex $x$ and a node $W$ of size $N$, let $h=\lceil4eDN/n+32\log n\rceil$. If $h\le D$, then, conditioned on $\Omega_0$, the probability that at least $h$ outneighbours of $x$ lie in $W$ is at most
\[
\binom Dh\Bigl(\frac{2N}n\Bigr)^h\le\Bigl(\frac{2eDN}{nh}\Bigr)^h\le2^{-h}\le n^{-32\log 2}.
\]
If $h>D$, then this event cannot occur. Taking a union bound over the at most $2n\cdot n\cdot n$ triples $(Z,x,W)$, we see that, with probability at least $1-2n^{3-32\log2}$ conditioned on $\Omega_0$, every node $W$ of size $N$ satisfies
\begin{equation}\label{eq:D2}
|E_Z(K[W])|\le N(h-1)\le\Lambda(N)\qquad\text{for all }Z\in\cZ.
\end{equation}

\noindent\textbf{Edges of one colour.}
A colour class is a matching with at most $M$ edges. For a node $W$ of size $N$ to contain $h'=\lceil8eMN^2/n^2+32\log n\rceil\le M$ of its edges, some $2h'$ specified vertices must lie in $W$, where $2h'\le2M\le n/4$. So, conditioned on $\Omega_0$, this has probability at most
\[
\binom M{h'}\Bigl(\frac{2N}n\Bigr)^{2h'}\le\Bigl(\frac{4eMN^2}{n^2h'}\Bigr)^{h'}\le n^{-32\log2}.
\]
Taking a union bound over the at most $n^2\cdot n$ pairs of a colour and a node, we see that, with probability at least $1-n^{3-32\log 2}$ conditioned on $\Omega_0$,
\begin{equation}\label{eq:D3}
|E_\gamma\cap E(K[W])|\le m(|W|)\qquad\text{for all colours }\gamma\text{ and nodes }W.
\end{equation}
Since $3n^{3-32\log2}<1$, we can fix a bijection $\sigma\in\Omega_0$ for which~\eqref{eq:D2} and~\eqref{eq:D3} hold. From now on, $\sigma$ is fixed.

\subsection{Decomposing the leaves of $\cT$}\label{sec:leaves}
Let $W$ be a leaf of $\cT$, of size $b\in[B,2B]$, and recall that $K[W]$ is rainbow as $\sigma\in \Omega_0$.
Take a Walecki decomposition of $K[W]$ into $b/2$ Hamilton paths, and relabel $W$ by an independent uniformly random bijection, so that each of these paths is a uniformly random Hamilton path of $K[W]$.
Fix one of these paths, $P$ say, and some $Z\in\cZ$. By~\eqref{eq:D2}, we have $|E_Z(K[W])|\le\Lambda(b)$, while $t(b)\ge 8e\Lambda(b)/b$ and $t(b)\le b/2$. So, by Lemma~\ref{lem:ham}, $\Prb(\ell_Z(P)\ge t(b))\le2^{-t(b)}\le n^{-2048\log 2}$.
Taking a union bound over the $n/2$ paths and the at most $2n$ sets $Z\in \cZ$, we can fix decompositions of all the leaves of $\cT$ into Hamilton paths $P$ with $\ell_Z(P)\le t(b)-1$ for all $Z\in\cZ$, where $b$ is the size of the leaf containing $P$.

\subsection{Combining the trees up the recursion}\label{sec:combining}
We now prove, by induction on the height of a node $W$ of $\cT$, that $E(K[W])$ decomposes into $|W|/2$ rainbow spanning trees of $K[W]$, such that each of these trees, $T$ say, has the following property. There is a path $W=W_0\supset W_1\supset\dots\supset W_k$ in $\cT$ from $W$ down to a leaf of $\cT$, such that
\begin{equation}\label{eq:load}
\ell_Z(T)\le\sum_{j=0}^k\bigl(t(|W_j|)-1\bigr)\le L(|W|)\qquad\text{for each } Z\in\cZ.
\end{equation}
Here, the second inequality in~\eqref{eq:load} always holds, by the properties of $\cT$ in Section~\ref{sec:hierarchy}, since
\[\textstyle\sum_j(1024D|W_j|/n+2048\log n+1)\le2560D|W|/n+2\log n\,(2048\log n+1).\]
For the leaves of $\cT$, the decompositions from Section~\ref{sec:leaves} satisfy~\eqref{eq:load}, with $k=0$.

Let $W$ be a node of $\cT$ of size $N$, with children $A$ and $A'$ of (even) sizes $a$ and $a'$, and suppose that the claim holds for $A$ and $A'$.
Split $A=A_1\cup A_2$ and $A'=A_1'\cup A_2'$ into halves, and let $H=E(A_1,A_1')\cup E(A_2,A_2')$ and $\barr H=E(A,A')\setminus H$. Then $d_H(v)=a/2$ for each $v\in A'$, and $d_{\barr H}(u)=a'/2$ for each $u\in A$.
We apply Lemma~\ref{lem:attach} to the $q=a/2$ trees on $A$, with $A'$ and $H$, using the parameters
\[M_{\rm loc}=m(N),\qquad K=2M,\qquad\Lambda=\Lambda(N)\qquad\text{and}\qquad t=t(N),\]
where $M_{\rm loc}$ takes the role of $M$ in Lemma~\ref{lem:attach}.
We now check the conditions of Lemma~\ref{lem:attach}.
\begin{itemize}
\item[\ref{att:mult}] For each colour $\gamma$, the edges counted by $a_\gamma+d_\gamma$ are distinct edges of colour $\gamma$ in $K[W]$, so that $a_\gamma+d_\gamma\le m(N)$ by~\eqref{eq:D3}.
\item[\ref{att:Z}] Each colour lies in at most $K=2M$ sets in $\cZ$ by~\ref{Z3}, and, for each $Z\in \cZ$, the number of edges of $H$ with colour in $Z$ is at most $|E_Z(K[W])|\le\Lambda(N)$ by~\eqref{eq:D2}.
\item[\ref{att:ineq}] For $v\in A'$, we have $\phi(H(v))\subseteq C(v)\subseteq Z_v$ by~\ref{Z4}. So, $r\le\max_i\ell_{Z_v}(T_i)\le L(a)\le a/1024$ by~\eqref{eq:load} for $A$ and by~\eqref{eq:consts}.
We have $q=a/2\ge (N-2)/4\ge N/5$. Moreover, $8r\le a/128\le N/200$, and $32(m(N)-1)\le1024MN/n+1024\log n\le N/1000$, using that $M/n\le2^{-48}$ and $N\ge B$. Hence $q\ge8r+32(M_{\rm loc}-1)$.
Next, $4e\Lambda(N)/q\le20e(16DN/n+32\log n)\le1024DN/n+2048\log n\le t(N)-1$.
Finally, $\log_2(16K)=\log_2(32M)\le5+\log_2n\le2048\log n\le t(N)-1$.
\end{itemize}
Thus, by Lemma~\ref{lem:attach}, we can extend the $a/2$ trees on $A$ to rainbow spanning trees of $K[W]$ using exactly the edges of $H$, so that, for each $Z\in \cZ$, the value $\ell_Z(T)$ of each tree $T$ increases by at most $t(N)-1$.
In the same way, we can extend the $a'/2$ trees on $A'$ using exactly the edges of $\barr H$, which is disjoint from $H$.
This gives $N/2$ rainbow spanning trees of $K[W]$ satisfying~\eqref{eq:load}, where, for a tree extending a tree on $A$, we extend its path in $\cT$ from $A$ by adding $W$ at the top.
Every edge of $K[W]$ is used exactly once, as it either lies inside a leaf of $\cT$, or lies in $E(U,U')$ for the two children $U$ and $U'$ of the lowest common ancestor in $\cT$ of the nodes containing its vertices.
This completes the induction. At the root of $\cT$, it gives the decomposition required in Theorem~\ref{thm:tests}, with $\ell_Z(T)\le L(n)=2560D+4097(\log n)^2$ for every tree $T$ and every $Z\in \cZ$. \qed
